\section{Conclusion}
\label{sec:conclusion}

This paper defines embodied AI evaluation deviation diagnosis as a software engineering problem and presents \sys, a manifest- and replay-based approach for classifying deviations, ranking engineering factors, and abstaining when evidence is insufficient. The empirical study shows that real projects contain diverse deviation symptoms and many under-specified factor reports, while the LeRobot+LIBERO evaluation shows that factor-directed replay can convert metric changes and failed executions into developer-facing diagnoses. Negative calibration, ablation, robustness, and cost results further show why replay evidence, conservative abstention, and targeted replay matter. The broader lesson is that embodied-AI benchmark results should be treated as outputs of an evaluation pipeline, not as policy-only measurements: a score change becomes actionable only when developers can connect it to manifest, episode-level, failure, and replay evidence.
